\scanpage{13}

\diagram{\begin{tikzcd}[ampersand replacement=\&,column sep=2em,row sep=2.2em]
1\arrow[r]\&Z\arrow[r,phantom,"\unsure{-|\to}" description]\&G\arrow[r,dash]\&\Gamma\arrow[r,dash]\&1\&[3em](E)\\
1\arrow[r,phantom,"\unsure{-|\to}" description]\&Z\arrow[u,"\unsure{2|\wr}"']\arrow[r,dash]\&G'\arrow[u]\arrow[r,dash]\&\Gamma'\arrow[u]\arrow[r,dash]\&1\&(E')
\end{tikzcd}}

If the extension $(E)$ \unsure{splits|splits}, the extension $E'$ too.

\noindent
\begin{tikzpicture}[baseline=(t.base)]
  \node[anchor=base west,inner sep=0pt] (t) at (0,0) {Every};
  \node[anchor=base west,inner sep=0pt] at (-.32,.60) {\inserted{Remark}};
  \draw (.64,.48) .. controls (.30,.32) and (-.43,.35) .. (-.39,-.06);
\end{tikzpicture}
split extension $G$ of a $\Gamma$ by a $Z$ arises by pullback from \deleted{\unsure{a}}the split extension
\[
\begin{tikzpicture}[baseline=(za.base)]
  \path[use as bounding box] (0,-.1) rectangle (12,.65);
  \node[anchor=base west,inner sep=0pt] (za) at (0,0) {$Z.\operatorname{Aut}(Z).$};
  % Detached two-lobed trace; writing versus surface trace remains uncertain.
  \node[inner sep=0pt] at (8.15,.40) {\unsure{\begin{tikzpicture}[baseline=0pt]
    \fill (-.006,.048) ellipse[x radius=.018cm,y radius=.032cm,rotate=24];
    \fill (.002,-.012) ellipse[x radius=.019cm,y radius=.037cm,rotate=12];
  \end{tikzpicture}}};
\end{tikzpicture}
\]

\figure{\begin{tikzpicture}
  \path[use as bounding box] (0,-.06) rectangle (12,.55);
  \draw (4.75,.28)--(7.12,.38);
  \draw (4.83,.08)--(6.96,.10);
  % Tiny detached point beneath the lower rule endpoint; possible surface trace.
  \node[inner sep=0pt] at (6.98,.015) {\unsure{\begin{tikzpicture}[baseline=0pt]
    \fill (0,0) ellipse[x radius=.009cm,y radius=.011cm];
  \end{tikzpicture}}};
  % Isolated point to the right of the divider; possible surface trace.
  \node[inner sep=0pt] at (10.02,.13) {\unsure{\begin{tikzpicture}[baseline=0pt]
    \fill (0,0) ellipse[x radius=.024cm,y radius=.021cm,rotate=-16];
  \end{tikzpicture}}};
\end{tikzpicture}}

4) \quad Sets of small cardinality

a) \uline{Cardinality 2}\qquad Sets of cardinality 2 $\longleftrightarrow$ \unsure{torsor|torsors} \unsure{over|under} $\mathbb{Z}_2$ or again, \uline{affine space \unsure{\deleted{under}|\deleted{over}|under|over}} $\deleted{\unsure{\mathbb{Z}|Z}}\mathbb{F}_2$ (or again, \deleted{\illegible{1}} affine space of dimension 1 over field $\mathbb{F}_2$ --- for the vector spaces of dimension 1 over $\mathbb{F}_2$ are canonically isomorphic \ldots)

\[
\begin{aligned}
&\begin{tikzpicture}[baseline=(bx.base)]
  \path[use as bounding box] (-.1,-.08) rectangle (11.3,1.45);
  \node[anchor=base west,inner sep=0pt] (bx) at (.35,.52)
    {$\operatorname{Aff}(1,\mathbb{F}_2)\simeq\operatorname{Aff}(\mathbb{Z}_2)\simeq\unsure{\mathfrak{S}|\mathcal{S}}_2$};
  % The drawn frame has gaps at the upper left and lower right.
  \draw (-.05,1.13)--(8.25,1.43)--(8.36,.43);
  \draw (0,.94)--(.09,0)--(8.30,.24);
  % Detached pair to the right of the box; possible surface trace.
  \node[inner sep=0pt] at (10.94,.67) {\unsure{\begin{tikzpicture}[baseline=0pt]
    \fill (.032,.032) ellipse[x radius=.020cm,y radius=.024cm,rotate=-22];
    \fill (-.020,-.038) ellipse[x radius=.024cm,y radius=.019cm,rotate=12];
  \end{tikzpicture}}};
\end{tikzpicture}\\[-.1em]
&\hspace{4.1cm}(\unsure{\mathcal{A}|A}_2=\{1\})
\end{aligned}
\]

b) \uline{Cardinality 3}
\[
\begin{aligned}
&\mathfrak{S}_3\simeq
\begin{gathered}[t]
  \operatorname{Aff}(1,\mathbb{F}_3)\\[-.15em]
  \rotatebox{90}{$\simeq$}\\[-.15em]
  \begin{aligned}
    \mathbb{F}_3^{+}&\mathbin{.}\mathbb{F}_3^{*}\\[-.4em]
    \rotatebox{90}{$\unsure{\simeq|=}$}&\quad\rotatebox{90}{$=$}\\[-.25em]
    \mathbb{Z}_3&\quad\mathbb{Z}_2
  \end{aligned}
\end{gathered}
\simeq\operatorname{Gl}(2,\mathbb{F}_2)\simeq
\begin{gathered}[t]
  \operatorname{GP}(1,\mathbb{F}_2)\\[-.15em]
  \rotatebox{90}{$=$}\\[-.15em]
  \begin{tikzpicture}[baseline=(qq.base)]
    \node[inner sep=0pt] (qq) at (0,0) {$\operatorname{Gl}(2,\mathbb{F}_2)/\mathbb{F}_2^{*}$};
    % Detached point below the quotient's left portion; possible surface trace.
    \node[inner sep=0pt] at (-1.25,-.58) {\unsure{\begin{tikzpicture}[baseline=0pt]
      \fill (0,0) ellipse[x radius=.010cm,y radius=.013cm];
    \end{tikzpicture}}};
  \end{tikzpicture}
\end{gathered}\\[.55em]
&\hspace{.3cm}\mathcal{A}_3\simeq\deleted{\mathbb{Z}}\mathbb{F}_3^{+}(\simeq\mathbb{Z}_3)
\end{aligned}
\]
% archive sheet number: 12
